\section{Related Work}
\label{sec:related}

\para{Long-term conversational memory.}
MemGPT~\cite{packer2024memgpt} introduces a hierarchical memory architecture inspired by OS virtual memory management, with working context and external archival storage. It supports persistent user facts across sessions but stores them as flat key-value pairs with no dependency structure, meaning a location change does not invalidate any downstream memories. MemoryBank~\cite{zhong2023memorybank} applies the Ebbinghaus forgetting curve to decay memory strength over time, providing a temporal signal for staleness detection. However, temporal decay treats all memories uniformly — a recently stated but stale preference is retained, while an older but still-valid fact decays. Unlike these systems, we use typed graph edges to propagate invalidation selectively based on semantic dependency, not just recency.

\para{Preference evolution and personalization benchmarks.}
\horizonbench~\cite{li2026horizonbench} is the most directly relevant prior work: it introduces a mental state graph with typed dependency edges and evaluates 25 frontier LLMs on preference evolution tracking over 6-month synthetic histories. The benchmark reveals that all models exhibit belief-update failure, selecting pre-evolution preferences as answers 25\%+ of the time. \horizonbench uses structured dependency edges to \textit{generate} evolving preferences, but evaluates models on raw context comprehension rather than on cascade propagation mechanisms. Our work takes the complementary approach: we implement and evaluate the cascade propagation mechanism itself on both real and structured data. PersonaMem-v2~\cite{jiang2025personamem} achieves state-of-the-art personalization accuracy (55\%) using a compact agentic memory updated across sessions, demonstrating that 2K-token persistent summaries outperform full-context retrieval. Building on this insight, our graph-based approach provides the structured update mechanism that agentic memory systems currently lack.

\para{Knowledge editing with ripple effects.}
\rippledits~\cite{cohen2023rippleedits} formalizes ripple effects in knowledge editing, defining six evaluation criteria for how a single factual edit should propagate through a knowledge graph. It shows that current knowledge editing methods fail to propagate changes consistently. Our work adapts this framework to user preference graphs: \locatedin edges correspond to compositional dependencies (editing a location should update location-dependent facts), and \conflictswith edges correspond to logical generalization (a new preference implies contradicting an old one). Unlike \rippledits, which focuses on encyclopedic world facts, we target personalized user preferences, where the conflict relationships are often implicit and require semantic inference rather than direct logical implication. ChainEdit~\cite{ChainEdit2025} combines knowledge graph rules with LLM reasoning for chain updates, closely paralleling our hybrid approach. Our work extends this to user preference graphs with a focus on the behavioral and semantic signals available in conversational data.

\para{Preference drift detection.}
PAMU~\cite{sun2025pamu} proposes a sliding-window exponential moving average approach for detecting preference drift across five dimensions (tone, length, emotion, density, formality) on the \locomo dataset. It provides the quantitative signal for \textit{when} drift has occurred, which complements our work on \textit{propagating} the resulting invalidation. Unlike PAMU, which detects changes within a single preference dimension at a time, our cascade approach handles cross-preference semantic conflicts (\eg quiet preferences invalidating bar recommendations). The two approaches are complementary: PAMU-style drift detection could provide the trigger signal for initiating our cascade.

\para{Benchmark datasets.}
\locomo~\cite{maharana2024locomo} provides the only publicly available large-scale real-world long-term dialogue dataset (50 dialogues, up to 300 turns each), which we use for structural cascade evaluation. PersistBench~\cite{persistbench2026} introduces safety-aware memory persistence benchmarks, motivating the importance of timely memory invalidation. Our work provides the mechanism that PersistBench identifies as missing: a principled way to automatically retire outdated memories as user contexts evolve.
